\documentclass[11pt,a4paper]{article}
\usepackage[margin=1in]{geometry}
\usepackage{amsmath,amssymb,amsthm}
\usepackage{algorithm}
\usepackage{algpseudocode}
\usepackage{booktabs}
\usepackage{hyperref}

\theoremstyle{definition}
\newtheorem{definition}{Definition}[section]
\newtheorem{theorem}{Theorem}[section]
\newtheorem{lemma}[theorem]{Lemma}
\newtheorem{remark}{Remark}[section]

\title{\textbf{Rigorous Mathematical Specification, Early-Training Stability Dynamics, and Functional Schedules of Learning Rate Warmup}}
\author{Team Scaibu \\ \small Email: \texttt{jaydeep.9664920749@gmail.com} \\ \small Architecture and Core Engine Team}
\date{October 2026}

\begin{document}
\maketitle

\begin{abstract}
This document presents the complete mathematical foundation and algorithmic specification of Learning Rate Warmup. In early deep neural network optimization—particularly in Transformers and deep ResNets—gradient variance is extremely high and second-moment statistics are uncalibrated. Initializing training with a large learning rate induces catastrophic gradient explosions or bad local saddle traps. Warmup gradually ramps the step size $\eta_t$ from an initial value $\eta_{\text{min}}$ to peak rate $\eta_{\text{base}}$ over $W$ steps. We formalize linear, cosine, and quadratic ramp dynamics, present the exact algorithmic pseudo-code matching the reference implementation, derive early variance stability theorems, analyze complexity, detail a numerical example, and discuss boundary conditions.
\end{abstract}

\section{Notation and Mathematical Preliminaries}

Let $t \in \mathbb{N}_{\ge 0}$ denote the discrete training step counter.

\begin{itemize}
    \item $t$: Current training step index ($t \ge 0$).
    \item $W \in \mathbb{N}_{\ge 1}$: Warmup duration (total warmup steps).
    \item $\eta_{\text{base}} > 0$: Target peak base learning rate ($\eta_{\max}$).
    \item $\eta_{\text{init}} \ge 0$: Initial starting learning rate at step $t=0$ (typically $0.0$).
    \item $p(t) \in [0, 1]$: Normalized warmup progress fraction, defined as $p(t) = \min(1.0, t / W)$.
    \item $\eta(t) \in [\eta_{\text{init}}, \eta_{\text{base}}]$: Scheduled learning rate at step $t$.
\end{itemize}

\begin{table}[h]
\centering
\caption{Summary of Mathematical Symbols for Learning Rate Warmup}
\begin{tabular}{lll}
\toprule
\textbf{Symbol} & \textbf{Type / Domain} & \textbf{Description} \\
\midrule
$t$ & $\mathbb{N}_{\ge 0}$ & Current step counter \\
$W$ & $\mathbb{N}_{\ge 1}$ & Warmup horizon ($1000$ to $5000$) \\
$\eta_{\text{base}}$ & $\mathbb{R}_{> 0}$ & Peak target learning rate \\
$\eta_{\text{init}}$ & $\mathbb{R}_{\ge 0}$ & Initial step-zero rate ($0.0$) \\
$p(t)$ & $[0, 1]$ & Progress ratio $t / W$ \\
$\eta(t)$ & $\mathbb{R}_{> 0}$ & Effective scheduled step size \\
\bottomrule
\end{tabular}
\end{table}

\section{Problem Definition and Core Concepts}

\subsection{Intuition and Motivation}
In post-LN Transformers and large-batch vision models, gradients during the first few hundred steps exhibit enormous variance because layer representations have not yet aligned. When combined with Adam's early second-moment estimates ($v_t \approx 0$), unconstrained step sizes produce massive parameter displacements ($\|\Delta \theta\| \gg 1$), knocking parameters out of well-behaved initialization basins. Ramping up the learning rate allows second-moment buffers to accurately accumulate gradient statistics under small, safe perturbations before exposing the network to peak learning rates.

\subsection{Formal Mathematical Definition}
\begin{definition}[Warmup Schedule Functions]
Given progress fraction $p = t / W$ for $0 \le t < W$:
\begin{enumerate}
    \item \textbf{Linear Warmup}:
    \begin{equation}
    \eta(t) = \eta_{\text{init}} + (\eta_{\text{base}} - \eta_{\text{init}}) \cdot p \label{eq:linear_warmup}
    \end{equation}
    \item \textbf{Cosine Warmup}:
    \begin{equation}
    \eta(t) = \eta_{\text{init}} + (\eta_{\text{base}} - \eta_{\text{init}}) \cdot \frac{1 - \cos(\pi p)}{2} \label{eq:cosine_warmup}
    \end{equation}
    \item \textbf{Quadratic Warmup}:
    \begin{equation}
    \eta(t) = \eta_{\text{init}} + (\eta_{\text{base}} - \eta_{\text{init}}) \cdot p^2 \label{eq:quad_warmup}
    \end{equation}
\end{enumerate}
For all steps $t \ge W$, $\eta(t) = \eta_{\text{base}}$ (or transitions to the subsequent decay schedule).
\end{definition}

\section{Algorithmic Specification}

The operational execution implemented in \texttt{NnAlgoLearningRateWarmup} is presented below.

\begin{algorithm}[H]
\caption{Learning Rate Warmup Step Computation}
\begin{algorithmic}[1]
\Require Current step $t \ge 0$, warmup steps $W \ge 1$, base rate $\eta_{\text{base}} > 0$
\Require Initial rate $\eta_{\text{init}} \ge 0$, schedule strategy $\in \{\text{linear}, \text{cosine}, \text{quadratic}\}$
\Ensure Current learning rate $\eta_t$, progress fraction $p$, warmup flag $\text{is\_warmup}$
\If{$t < 0$ \textbf{or} $W < 1$ \textbf{or} $\eta_{\text{base}} \le 0$ \textbf{or} $\eta_{\text{init}} < 0$}
    \State \textbf{throw} PreconditionFailureException(``Invalid schedule arguments'')
\EndIf
\If{$t \ge W$}
    \State \Return $\{\text{learning\_rate}: \eta_{\text{base}}, \text{progress\_fraction}: 1.0, \text{is\_warmup}: \textbf{false}\}$
\EndIf
\State $p \gets \text{float}(t) / \text{float}(W)$
\If{$\text{strategy} = \text{linear}$}
    \State $\eta_t \gets \eta_{\text{init}} + (\eta_{\text{base}} - \eta_{\text{init}}) \cdot p$
\ElsIf{$\text{strategy} = \text{cosine}$}
    \State $\text{factor} \gets 0.5 \cdot (1.0 - \cos(\pi \cdot p))$
    \State $\eta_t \gets \eta_{\text{init}} + (\eta_{\text{base}} - \eta_{\text{init}}) \cdot \text{factor}$
\ElsIf{$\text{strategy} = \text{quadratic}$}
    \State $\eta_t \gets \eta_{\text{init}} + (\eta_{\text{base}} - \eta_{\text{init}}) \cdot (p^2)$
\Else
    \State \textbf{throw} PreconditionFailureException(``Unrecognized strategy'')
\EndIf
\State \Return $\{\text{learning\_rate}: \eta_t, \text{progress\_fraction}: p, \text{is\_warmup}: \textbf{true}\}$
\end{algorithmic}
\end{algorithm}

\section{Theoretical Guarantees and Variance Analysis}

\begin{theorem}[Second-Moment Calibration Under Linear Warmup (Liu et al., 2019)]
Let gradients be drawn with variance $\sigma^2$. The variance of the Adam step $\Delta \theta_t = \frac{\eta_t \hat{m}_t}{\sqrt{\hat{v}_t}}$ is unbounded for $t < \frac{2}{1 - \beta_2}$ without warmup. Under linear warmup $\eta_t = \eta_{\text{base}} \frac{t}{W}$ with $W \ge \frac{5}{1 - \beta_2}$, the variance of the parameter update is bounded for all $t \ge 1$:
\begin{equation}
\operatorname{Var}(\Delta \theta_t) \le C \eta_{\text{base}}^2
\end{equation}
for a universal constant $C < \infty$.
\end{theorem}

\begin{proof}
For small $t$, the effective degrees of freedom $\rho_t = 2 / (1 - \beta_2) - 1$ of the scaled $t$-distribution in second-moment estimation is non-convergent when $\rho_t \le 4$. Scaling $\eta_t \propto t$ directly dampens the numerator, canceling the divergence in $\mathbb{E}[1 / \hat{v}_t]$.
\end{proof}

\subsection{Limitations}
\begin{itemize}
    \item \textbf{Heuristic Duration}: Setting $W$ too small fails to stabilize moments; setting $W$ excessively long wastes compute budget at sub-optimal learning rates.
\end{itemize}

\section{Complexity Analysis}

\subsection{Time and Space Complexity}
\begin{itemize}
    \item \textbf{Time}: Exactly $\mathcal{O}(1)$ scalar arithmetic operations (one division, one trigonometric or polynomial evaluation).
    \item \textbf{Space}: $\mathcal{O}(1)$ auxiliary memory.
\end{itemize}

\section{Worked Numerical Example}

Consider $W = 1000$, $\eta_{\text{base}} = 0.001$, $\eta_{\text{init}} = 0.0$, at step $t = 250$.

\subsection{Calculations}
\begin{enumerate}
    \item Progress fraction: $p = 250 / 1000 = 0.25$.
    \item \textbf{Linear}:
    \begin{equation}
    \eta_{\text{lin}} = 0.0 + (0.001 - 0.0) \times 0.25 = 0.000250
    \end{equation}
    \item \textbf{Cosine}:
    \begin{equation}
    \text{factor} = 0.5(1 - \cos(0.25 \pi)) = 0.5(1 - 0.7071068) = 0.5(0.2928932) \approx 0.1464466
    \end{equation}
    \begin{equation}
    \eta_{\text{cos}} = 0.001 \times 0.1464466 \approx 0.0001464
    \end{equation}
    \item \textbf{Quadratic}:
    \begin{equation}
    \eta_{\text{quad}} = 0.001 \times (0.25^2) = 0.001 \times 0.0625 = 0.0000625
    \end{equation}
\end{enumerate}

\section{Numerical Considerations and Edge Cases}

\begin{itemize}
    \item \textbf{Step Zero Evaluation}: At $t=0$, $\eta(0) = \eta_{\text{init}}$. Setting $\eta_{\text{init}} = 0$ prevents parameter change on step 0 while allowing momentum buffers to initialize on the first gradient.
    \item \textbf{Integer Division}: In Python 2 / C++ implementations, $t / W$ must be cast to floating point to avoid truncation to zero.
\end{itemize}

\begin{thebibliography}{9}
\bibitem{goyal2017} P.~Goyal et al., ``Accurate, large minibatch SGD: Training ImageNet in 1 hour,'' \emph{arXiv preprint arXiv:1706.02677}, 2017.
\bibitem{vaswani2017} A.~Vaswani et al., ``Attention is all you need,'' in \emph{Advances in Neural Information Processing Systems (NeurIPS)}, pp.~5998--6008, 2017.
\bibitem{liu2019} L.~Liu et al., ``On the variance of the adaptive learning rate and beyond,'' in \emph{International Conference on Learning Representations (ICLR)}, 2020.
\end{thebibliography}

\end{document}
